\section{失败、反例和适用范围}
\label{sec:03_08}

第\ref{sec:03_07}节的结果同时包含收益差距、停止时尚未闭合的最优性间隙，以及未取得可验收输出的运行。这些现象回答不同的问题：收益差距说明一个可行选择还有多大改进空间，最优性间隙说明证明进行到什么程度，执行异常则说明某次调用没有完成预定的输出与验收过程。只有把这些含义分开，才能确定一项不利结果究竟限制了哪一个主张。

\subsection{可行性、最优性与运行终止的区别}

沿用前文的有限候选集合$E$、资源集合$V$及关联矩阵$A$。候选$e\in E$的选择量为$x_e\in\{0,1\}$，固定收益为$w_e>0$；$A_{ve}=1$表示任务$e$需要资源$v$的一份容量，否则为零。当前组合实验均取$b_v=1$。对一个已经取得的二元输出$\widehat x$，可行性要求每个资源的占用不超过一份，收益则按原输入重新求和。记其核验收益为$h=\sum_{e\in E}w_e\widehat x_e$，相同问题的最优收益为$z^\star$，二者均无量纲。

合法次优解要求已经知道$\widehat x$可行且$h<z^\star$。即使不知道精确的$z^\star$，只要另有一个核验可行解的收益严格高于$h$，也足以证明该输出次优。最优性未证则只表示现有证据不能判定$h=z^\star$：该输出可能已经最优，也可能仍可改进。因而，缺少最优性证书不能直接写成已经证明次优。

表\ref{s0308:status_meanings}将解的性质与运行状态分别列出。这些项目并非互斥分类。例如，一次资源保护退出可以留下待核验候选；一次正常结束的启发式调用可以返回合法次优解；一次限时结束的SCIP调用也可以留下可行解。

\begin{table}[tbp]
 \centering
 \small
 \caption{结果性质与运行状态的证据含义}
 \label{s0308:status_meanings}
 \begin{tabularx}{\linewidth}{@{}p{0.15\linewidth}XX@{}}
  \toprule
  项目 & 可以据此确认的内容 & 不能直接推出的内容\\
  \midrule
  合法次优解 & 原约束核验通过，且存在严格更好的可行收益或已知最优值 & 程序执行失败，或该解没有应用价值\\
  最优性未证 & 尚无足够证据确认当前收益达到全局最优 & 当前解必然次优，或原问题无解\\
  工程异常 & 某次调用在接口、原生执行或验收环节异常终止 & 原论文算法一般失效，或异常根因已经定位\\
  资源保护 & 外部资源管理条件触发了停止 & 算法正常完成，或已证明所需内存的严格峰值\\
  未完成或未验收 & 预定计算或输出验收尚未完成 & 未记录的输出可自动视为有效，或其目标值等于零\\
  \bottomrule
 \end{tabularx}
\end{table}

这一划分还排除了把求解失败解释为模型不可行的做法。对本章的基本选择模型，全零向量始终满足单位容量约束，式\eqref{s0203:zero_feasible}已经证明其可行域非空。某次程序没有返回可验收结果，不能推翻这一模型性质。全零选择的收益确实为零，但失败记录并不是程序返回并核验了全零选择的证据。

\subsection{收益反例限定的算法主张}

十二任务算例提供了一个可以完整核查的反例。枚举覆盖全部$4096$个二元选择，其中$62$个可行，唯一最优选择为$\{\mathrm{T02},\mathrm{T03},\mathrm{T04},\mathrm{T12}\}$，收益为$42$。Streaming返回$\{\mathrm{T01},\mathrm{T05}\}$，收益为$39$，且两种选择均将每项资源恰好占用一次。于是，该例同时满足
\begin{equation}
 \ell_v(x^{\mathrm S})=\ell_v(x^\star)=1,
 \quad\forall v\in V,
 \qquad
 z(x^{\mathrm S})=39<42=z(x^\star).
 \label{s0308:full_load_counterexample}
\end{equation}
这里$x^{\mathrm S}$与$x^\star$分别为第\ref{sec:02_07}节定义的Streaming选择和最优选择，$\ell_v$为资源$v$的占用份数，$z$为无量纲收益函数。相同的满负载对应不同收益，说明资源使用充分并不是收益最优的充分条件。Streaming输出已经不能直接追加任何非空资源任务，但撤销原选择后，仍能换入收益更高的组合。

该反例否定了Streaming在所有输入上都取得最优解，或都保持最优收益$95\%$的主张，因为$39/42\approx0.9286$。它没有否定第\ref{sec:03_02}节所述的、具有明确前提的近似保证\cite{ReinstadtlerFPU2025}，也没有否定这次输出的可行性。研究结论因此应落在具体的收益与成本取舍上，而不是把一次失优扩大为方法没有作用。

压力测试进一步限制了统一的高质量主张。59个条件的236次输出均通过容量和收益核验，但Streaming在S26的相对可行值比约为$0.6015$，Greedy在S28为$0.7500$，ILS1000在S16约为$0.8604$。这里的分母都是对应SCIP可行收益，未统一证明为最优值。这些反例已足以说明：三种方法都不能在所测全部条件上保证达到该参照收益的$95\%$。可行性全部通过与质量判据未全部通过并不矛盾。

时间与质量也须在同一输出上配对判断。例如，第\ref{sec:03_07}节的基准配置中，ILS1000平均质量比约为$1.0128$、平均时间比约为$2.1058$，但同时满足质量比不小于$0.95$和时间比不小于$2$的种子只有一个。两个均值分别过线，不能替代逐种子的联合判定。三个诊断条件还显示，增大停止阈值未必严格增加收益：C06的ILS1000与ILS5000均取得$770$，调用时间却分别约为$1.31\,\mathrm{s}$和$5.96\,\mathrm{s}$。这限制了更大阈值必然改善结果的说法；由于两次调用独立启动，也不能将其未经核验地解释为同一搜索轨迹的两个截点。

\subsection{原作者完整流程的五次失败及统计分母}

原作者完整流程在七个既有输入上各运行十个固定算法种子，共有70次预定调用，其中65次取得通过原始权重与容量核验的输出，5次没有可验收结果。该流程包括至多10轮归约、$S,\mathrm{pin}$贪心、ILS15和展开\cite{GrossmannJRS2026}。表\ref{s0308:failure_details}保留五次失败的输入、种子及可确认的终止现象。

\begin{table}[tbp]
 \centering
 \small
 \caption{原作者完整流程中未取得可验收结果的五次运行}
 \label{s0308:failure_details}
 \begin{tabularx}{\linewidth}{@{}llXX@{}}
  \toprule
  输入 & 种子 & 已观察现象 & 解释边界\\
  \midrule
  \texttt{ibm02} & 7 & 可行性与收益一致性的联合断言失败 & 日志不能区分哪一项失败，也不能确认二者是否同时失败\\
  \texttt{us04} & 1 & 原生调用抛出Windows异常\texttt{0xc00000ff} & 未定位到具体内存、栈或其他内部原因\\
  \texttt{us04} & 3 & RAM保护触发退出 & 无可验收输出；保护退出不表示正常搜索结束\\
  \texttt{us04} & 7 & 原生调用抛出Windows异常\texttt{0xc00000ff} & 与种子1具有相同异常码，不据此推定相同内部根因\\
  \texttt{us04} & 9 & RAM保护触发退出 & 无可验收输出；不能将收益记为零\\
  \bottomrule
 \end{tabularx}
\end{table}

\texttt{ibm02}的联合断言尤其需要保留逻辑含义。令$P$表示返回选择满足原容量约束，$Q$表示按原权重重算的收益与程序报告一致，二者均为真假命题。日志确认的是
\begin{equation}
 \neg(P\land Q)=\neg P\lor\neg Q.
 \label{s0308:assertion_logic}
\end{equation}
它说明至少一项检查未通过，却不蕴含单独的$\neg P$。因此，这条记录不能写成已经发现一个违反容量的解；收益不一致同样可能触发该联合断言。现有日志也不足以把失败定位到某一条归约或展开规则。

为描述这些失败对统计的影响，令$p$标识一个公开输入，$\mathcal S_p$为其十个计划种子的集合，$\mathcal V_p\subseteq\mathcal S_p$为取得有效输出的种子集合。集合元素是无量纲运行标识。定义观察失败比例
\begin{equation}
 f_p=\frac{|\mathcal S_p\setminus\mathcal V_p|}{|\mathcal S_p|}
 =1-\frac{|\mathcal V_p|}{10},
 \qquad f_p\in[0,1].
 \label{s0308:failure_fraction}
\end{equation}
$f_p$描述当前固定种子和执行环境下的记录比例，不是已经估计出的跨平台失败概率。有效子集均值仍按式\eqref{s0307:author_quality_mean}计算，并须与$f_p$并列。表\ref{s0308:conditional_statistics}列出两个存在失败的输入。

\begin{table}[tbp]
 \centering
 \small
 \caption{存在失败时的有效子集收益与计划覆盖}
 \label{s0308:conditional_statistics}
 \begin{tabular}{@{}lrrrrl@{}}
  \toprule
  输入 & 有效/计划 & 失败比例 & 有效收益均值 & 有效收益范围 & 全十种子均值\\
  \midrule
  \texttt{ibm02} & $9/10$ & $10\%$ & $379719.3333$ & $379495$--$379987$ & 未定义\\
  \texttt{us04} & $6/10$ & $40\%$ & $978$ & $933$--$1065$ & 未定义\\
  \bottomrule
 \end{tabular}
\end{table}

表中收益均为无量纲原权重之和。其余五个输入均为十次有效，才具备计算全十种子收益均值的完整输出。整体$5/70$约为$7.14\%$，但整体比例不能掩盖失败集中在两个输入上的事实。删除失败后得到的均值只描述有效输出子集；没有证据表明缺失与所得收益无关，故不能把它当作全部计划运行的收益均值。有效子集的平均调用时间也不包含失败尝试的全部投入，不能据此代表为完成整个计划付出的总成本。

这些失败是当前运行环境与移植执行中的观察。公开源码与本地适配范围的对应材料在附册中保留，优化内核未被重写；但源码来源本身并不能证明适配、执行及验收的每一环节都正确。五次失败限定了本次完整流程的完成率及可报告统计，不足以断言原论文算法错误，也不足以断言它在其他环境或输入上一般失效。

\subsection{已经核验的浮点边界与尚未定位的原因}

第\ref{sec:03_03}节将理想实数优先级与实际机器表达式分开说明。对\texttt{us04}的独立表达式核验已经确认：在163项任务中，118项任务的资源度连乘在所检查的IEEE~754双精度表达式中溢出为正无穷，相应机器优先级成为零。这是已取得证据的数值表达式边界，不能笼统写成浮点异常是否存在尚不清楚。

其组合含义是，式\eqref{s0303:unit_priority}中严格为正、可能彼此不同的理想优先级，在该机器表达式下有一批被合并为零分。此时，数学分数所规定的严格大小关系未必还能由机器分数区分，具体访问顺序需要结合并列处理规则解释。该事实不意味着这些任务的原始收益为零，也不意味着它们违反资源容量。

这次核验检查的是保留适配器中的双精度算术表达式，没有重新执行原生贪心过程，也没有修改公开优化内核。因此，已核验的零分合并不能直接解释某次完整运行的收益差、异常退出或RAM保护。完整流程使用的$S,\mathrm{pin}$静态分数见式\eqref{s0304:static_pin_score}，与第\ref{sec:03_03}节的度乘积优先级不同，更不能仅因输入同为\texttt{us04}就把两类现象建立为因果关系。

原生内核其他浮点运算的行为、机器整数中间量是否溢出，以及哈希容器遍历是否造成跨运行差异，仍须各自的执行证据才能判定。尤其是ILS接受规则中的整数乘积，其数学解释以中间运算没有溢出为条件；这一条件尚不能改写为本批运行已经发生整数溢出。同样，固定随机种子说明随机输入受到控制，但仅凭种子值不能证明所有平台上的容器遍历和搜索轨迹一致。

\subsection{停止协议、相关样本与资源测量的边界}

SCIP的$5\%$间隙停止与启发式停止承担不同任务。前者在有效上下界和正可行收益的条件下限制尚未排除的最优收益范围；Streaming结束表示完成既定扫描与逆序构造；ILS停止表示连续未刷新最好收益的计数超过阈值。ILS15、ILS1000和ILS5000中的数值都是这一计数阈值，改进后重置，不是固定总迭代次数，也不是墙钟时间预算。

因此，前节的时间比可以说明相应调用按各自规则结束时的耗时关系，却不能直接构成等预算速度排名。初始记录、压力测试和完整流程还存在热启动、计时边界及CPU并发差异。续跑继承的可行解具有先前获取成本；历史首次正收益记录也不能被解释为重新实施了统一时限实验。求解器运行中的目标和界只是报告快照，在缺少对应见证核验时不能作为已完成决策。

样本数量也须对应实际生成过程。92项实验条件中包含共享拓扑、改变权重以及具有嵌套关系的派生条件；七个完整流程输入已经包含其中。算法种子重复观察同一输入上的搜索变化，结构种子则参与输入生成，两者并非同一种重复。十二任务算例的21次重复测时没有增加独立实例数。由这些相关记录计算的均值、范围和联合支持数具有描述意义，不能把全部运行视为相互独立的样本，进而补出没有实施的显著性结论。

内存记录还有采样边界。设一次进程观察区间为$[0,T]$，其中$T>0$以$\mathrm{s}$计；$M(t)\geq0$表示时刻$t$的驻留集大小，即RSS，统一以字节计。若仅在$k$个时刻$t_1,\ldots,t_k\in[0,T]$采样，$k$为正整数，则在采样值准确且统计口径一致的条件下，
\begin{equation}
 M_{\mathrm{sample}}=\max_{1\leq i\leq k}M(t_i)
 \leq
 M_{\mathrm{sup}}=\sup_{t\in[0,T]}M(t).
 \label{s0308:memory_sampling}
\end{equation}
$M_{\mathrm{sample}}$是采样最大RSS，$M_{\mathrm{sup}}$是整个观察区间的RSS上确界，二者单位相同。采样间隙内可能存在未捕获的增长，因此不能将采样最大值标为严格峰值。RAM保护记录说明保护条件被触发，并不单独给出算法完成所需的最小内存；RSS与其他内存口径也不能混用。

\subsection{从组合选择到巡检执行的适用范围}

本章实验直接适用于候选任务及关联已固定、收益为正且可加、每个关联需求占用一份资源、全部容量为一的选择问题。对任意二元向量$x\in\{0,1\}^{|E|}$，有
\begin{equation}
 \left[\sum_{e\in E}A_{ve}x_e\leq1,
       \ \forall v\in V\right]
 \quad\Longleftrightarrow\quad
 \left[x_e+x_f\leq1,
       \ \forall e\ne f:\ R_e\cap R_f\ne\varnothing\right].
 \label{s0308:unit_capacity_boundary}
\end{equation}
这里$e,f\in E$为不同任务，$R_e,R_f$为各自需要的资源子集。左侧逐资源限制占用，右侧逐冲突任务对限制共选；等价性依赖二元域及单位容量。较大的超边秩只表示单项任务关联更多资源，并不破坏这一等价性。因此，当前实验没有仅凭高秩证明不可约高阶耦合，也没有通过表示之间的受控比较证明超图形式必然更快。整数可行域等价还不意味着连续松弛、存储成本和实际求解时间相同，相关区别见第\ref{sec:02_04}节。

这些边界同样限制向无人集群巡检的外推。纯组合选择没有给出连续轨迹、观测姿态、电池能量变化、车辆与无人机会合时序，以及设备服务后的再次使用过程。共享资源检查通过，只说明当前静态占用定义下没有冲突。要成为可执行巡检方案，还必须为同一选择提供共同满足物理与时序条件的执行见证；逐个任务可行并不保证任务之间能够同步执行。第\ref{sec:01_03}节的投影关系也以资源约束确实是执行过程的必要条件为前提，不能自动套用于任意静态设备占用定义。

本章记录的秒数是主机上的计算耗时，不是无人机飞行历时或现场测量。公开MIS输入的复用也不构成HypOp神经求解器的性能实验，因为该求解器没有运行。由本章收益与时间数字，既不能证明工程联合模型更优，也不能推出某个区域可以实现完整覆盖。

区域模型已有有限转录的构建与运行记录；没有取得新结果不能改写为没有实现。现有区域记录中，LP求解错误属于较早的绑定运行；截至2026年11月10日03:56 UTC的后续观察，最新运行因资源保护退出，留下的四塔候选仍需重新验收，尚无新验收的完整区域覆盖方案。候选目标数值相同不能代替执行见证一致，也不能继承其他版本的验收结论。这一状态说明当前证据尚未给出完整覆盖方案，并不证明原联合问题无解；其目标允许部分完成，亦不能悄然改写为全部杆塔必选。

由此进入第4章时，需要补充的是选择之外的物理对象与状态关系：先定义地理位置、观测要求和设备，再逐项建立运动、成像、安全与能量条件。只有这些条件与后续分配、航次和同步约束共同成立，组合层的选择才可能获得巡检层的执行含义。
